%!TeX root=../quali.tex
\chapter{Introduction}
\label{ch:introduction}

High-performance computing (HPC) has become essential to modern science, and its capability has grown by drawing more power, packing more hardware, and cycling through that hardware ever faster \parencite{portegieszwart2020Ecological, allen2022Huge}. The environmental cost of this growth has traditionally been equated with operational energy consumption \parencite{koomey2011Implications}, yet operational energy is only one facet of the footprint. Manufacturing the processors, memory, and interconnects that populate a machine emits carbon long before the system is switched on, and this embodied share is rising as fabrication grows more complex \parencite{gupta2021Chasing, pirson2023Environmental}. Cooling those systems and generating the electricity that feeds them consumes freshwater at scales large enough to matter regionally \parencite{mytton2021Data}. The semiconductor supply chain draws on scarce materials and discharges hazardous waste \parencite{sandhu2025Semiconductor}, while the shortening replacement cycles that keep systems competitive turn working hardware into electronic waste. Across its full lifecycle, then, an HPC system carries costs in carbon, water, and materials and leaves electronic waste behind, well beyond the operational energy it draws.

Visibility into this broader footprint has only recently begun to grow, and the established ways of evaluating and operating HPC systems still capture only part of it. The Green500 ranks systems by performance per watt under a single benchmark \parencite{feng2007Green500}, reducing sustainability to an efficiency ratio measured on Linpack. Job schedulers decide when and where computation runs, yet they have conventionally been designed to prioritize performance above all else \parencite{kocot2023EnergyAware, czarnul2019EnergyAware}, leaving environmental cost outside the optimization objective. Procurement and hardware-replacement decisions weigh performance and acquisition cost while treating lifecycle impact as an afterthought \parencite{li2023Sustainable}. The research concerning these overlooked aspects has started to gain traction only lately. At the same time, current strategies target just one piece of a challenge that is not yet fully understood. Even the most closely monitored indicator shows no net progress, because rising demand continually cancels out any benefits achieved \parencite{york2022Rebound}.

To translate this rising attention into well-defined research needs, the project begins with a systematic mapping study \parencite{rosa2026Environmental}. Classifying 62 reports from sixteen years of research across environmental impacts, lifecycle stages, system loci, and management levers, the study finds the literature concentrated within a narrow path. In particular, existing papers focus mostly on carbon among the impacts, operations among the lifecycle stages, facility infrastructure among the system loci, and measurement among the management levers. Meanwhile, other topics appear only sporadically and are analyzed in isolation, such as non-cooling water usage, biodiversity, and end-of-life processing. To address these gaps, the rest of this work divides into two branches (\Cref{fig:plan-branches}).

\begin{figure}
\centering
\includegraphics[width=\textwidth]{branches}
\caption{The two branches of this work and the steps of each. Steps are ordered by dependency.}
\label{fig:plan-branches}
\end{figure}

The first branch focuses on the operation of HPC systems. To evaluate different management strategies without needing access to a production system, we built a simulation testbed. The testbed has so far introduced footprint models into existing simulators to enable the evaluation of operational carbon, operational water, and embodied footprint. A first result along this line is Greenfilling, a small change to EASY backfilling \parencite{skovira1996EASY} that holds back opportunistic job starts when the targeted grid intensity, carbon or water, runs above its recent average. Evaluated across three national grids and two production machines, it traded certain delay for uncertain environmental return. This result sets a lower bound on how sophisticated such an intervention has to be. A job-shifting variant was also developed. It prices delaying each job against the grid intensities it would meet, and an offline analysis of the same workloads and signals places its ceiling near 7\% of operational carbon, two orders of magnitude above what Greenfilling reached. We plan to extend this testbed by including user-facing accounting, platform controls, and lifecycle-aware policies.

The second branch focuses on the overall impacts of HPC systems. The Green500 ranks systems by performance per watt, but this is only one facet of the broader impacts of using those platforms. Two systems can rank closely in the Green500 while drawing very different absolute power, sitting on grids whose water and carbon intensity differ by orders of magnitude, carrying different embodied impacts, and consuming water at rates that do not reflect regional scarcity. By applying a sufficiency perspective to Green500 data, we intend to examine cases where efficiency improves while absolute burden grows, where embodied costs dominate but stay invisible, and where regional context changes which system is actually less harmful.

The remainder of this document is organized as follows. \Cref{ch:mapping} reviews the environmental footprint of HPC across its lifecycle, then presents the systematic mapping study and the classification behind the gap this work addresses. \Cref{ch:preliminary} reports the preliminary results. \Cref{ch:plan} sets out what remains. It adds user modeling, platform controls, and lifecycle-aware policies to the testbed built in the first branch, and designs a sufficiency-aware evaluation that uses the Green500 as a reference point, widening the ranking to the impacts efficiency omits and reaching toward the absolute burden a per-watt ratio hides, as far as public data and what-if scenarios permit. \Cref{ch:schedule} closes with the schedule and the expected contributions of the completed work.
